% !TeX root = ../main.tex
\newcounter{paperalgorithm}
\newenvironment{paperalgorithm}[2]{%
  \par\Needspace{12\baselineskip}\refstepcounter{paperalgorithm}\label{#2}%
  \begingroup\small\setstretch{1.05}\setlength{\parskip}{3pt}%
  \noindent\rule{\linewidth}{0.6pt}\par\noindent\textbf{算法\thepaperalgorithm\quad #1}\par
}{\par\noindent\rule{\linewidth}{0.6pt}\par\endgroup}

\begin{paperalgorithm}{窗口求解与逐时复核流程}{alg:season-workflow}
\noindent\textbf{输入：}场景配置、规划输入、连续窗口逐时数据与求解容差。\par
\noindent\textbf{输出：}投资及逐时运行方案，求解状态、间隙与校核记录。
\begin{enumerate}[label=\arabic*.,leftmargin=1.8em,itemsep=2pt,topsep=3pt,parsep=0pt]
\item 核对窗口、步长、单位与各模块启用状态。
\item 按第2章约束装配源网荷储电模型，联立守恒、容量与热状态条件并设初末边界。
\item 乌鲁木齐联合求解投资与运行；固原、武汉冻结投资模式后求解全季运行。
\item 按碳帽或 $\varepsilon$ 档位重复求解并逐档保存。
\item 记录目标值、上下界、间隙、时间与可行性状态。
\item 以逐时输出复算守恒、状态、容量、资源与失供核对。
\item 导出可核结果并标注求解身份与逐时调度。
\end{enumerate}
\end{paperalgorithm}
